\chapter{الهندسة الإهليليجية (الناقصية)}

\section[مقدمة]{مقدمة \cite{ref6}}
ان اول من اشار الى هذه الهندسة هو العالم الالماني ريمان في عام ١٨٥٤ في رسالة الدكتوراه حيث اخذ نقيض بديهية بليفر " وهي احدى مكافئات البديهية الخامسة لاقليدس " وكما يلي :
البديهية المميزة للهندسة الاهليجية:
" لا يمكن رسم أي مواز لخط معلوم من نقطة خارجة عنه".
وفي هذه الهندسة ميز ريمان بين فكرة المستقيم الغير منتهي وفكرة المستقيم الغير محدود .
(حيث سابقا لا يوجد فرق) وبذلك اخذ الخطوط اقواس لدوائر عظيمة على سطح كرة نصف قطرها الثابت a . فان الخط منتهي بالطول ، حيث له محيط معلوم ولكن غير محدود في المفهوم الذي نستمر فيه حول الكرة دون توقف.

اذا الكرة هي نموذج لهذه الهندسة.
نقاط هذا النموذج هي نقاط على سطح الكرة.
وان المستقيمات لهذا النموذج هي دوائر عظيمة. بما ان الدوائر العظمى تتقاطع دائما ، فان المستقيمين يتقاطعان دائما ، ولذلك لا توجد مستقيمات موازية للمستقيم.
بما ان توجد عدد غير منته من دوائر عظمى تمر بنقطتين متقابلتين بالقطر، فان النقطتين لا تعينان دائما مستقيما واحدا فقط. وبذلك فان بديهيات الوقوع غير متحققة. والتحول من الهندسة الاقليدية الى الهندسة الاهليجية لا يتم بسهولة كما يتم مع التحول من الهندسة الاقليدية الى الهندسة الهذلولية.

\subsection*{الهندسة الإهليليجية (الناقصية)}
ان الهندسة الاهليجية قائمة على اساس النقيض الثاني لبديهية اقليدس الخامسة وهي لا يوجد أي مواز للمستقيم أي ان جميع المستقيمات متقاطعة. ان النموذج الذي افترض لتحقيق هذه الهندسة هو سطح كرة مع اعتبار ان المستقيمات هي أجزاء من دوائر عظمى (الدائرة العظمى بالنسبة للكرة هي الدائرة التي يكون نصف قطرها يساوي نصف قطر الكرة) .

نفرض ان معادلة الكرة $x^2 + y^2 + z^2 = 1$\\
مركز هذه الكرة هو نقطة الأصل (0,0,0)\\
ونصف قطرها = 1.\\
إن أي مستوي يمر بالمركز يقطع سطح الكرة بدائرة عظمى.

يوجد نوعان من الهندسة الاهليجية:
\begin{enumerate}
	\item  الهندسة الإهليليجية الاحادية (Single E.G) والتي يكون فيها أي مستقيمين يتقاطعان في نقطة واحدة فقط. وان المستوى الإهليليجي يتميز بان سطحه له جهة واحدة . وكمثال بسيط على ذلك ورقة موفيس، ونموذج اخر عليها هو نصف كرة فقط .
	
	\item  الهندسة الإهليليجية الثنائية (Double E.G) والتي يكون فيها أي مستقيمين يتقاطعان في نقطتين . وان المستوى الإهليليجي يتميز بان سطحه له جهتين . مثلا الكرة التي فيها جهتين
\end{enumerate}


\begin{theorem}
	أي مستقيمين في الهندسة الإهليليجية المزدوجة يتقاطعان في نقطتين ويحيطان منطقة لمساحة منتهية A حيث ان $A \neq 0$ . او الخط المار بنقطتين متقابلتين ليس وحيد .
\end{theorem}

سنوضح عددا من المبرهنات المشتركة للهندسة الاحادية والمزدوجة .

\begin{theorem}
	العمودان على نفس المستقيم يتقاطعان في نقطة . مثلا تقاطع خطوط الطول والعرض على سطح الكرة الارضية.
	نستنتج هذه المبرهنة من البديهية المميزة لهذه الهندسة ، حيث ان الخطين في المستوي يتقاطعان دائما.
\end{theorem}

\begin{theorem}
	الاعمدة على كل النقاط لمستقيم تلتقي في نقطة تدعى قطب الخط، وبالعكس كل مستقيم يمر بالقطب يكون عموديا على المستقيم . مثل خطوط الطول هي نماذج للخطوط العمودية على خط الاستواء ، وان كل هذه الخطوط تتقاطع في القطبين الشمالي والجنوبي .
	المسافة q من قطب مستقيم للمستقيم ، هي المسافة القطبية للمستقيم .
\end{theorem}

\begin{theorem}
	في أي مثلث ABC , الذي فيه $\angle A = \angle B = \angle C = 90$ تكون اصغر من، تساوي او اكبر من 90 نسبة الى ان القطعة BC اصغر من ، تساوي او اكبر من المسافة القطبية q .
\end{theorem}

\begin{theorem}
	مجموع زوايا المثلث هو اكبر من 180 .
\end{theorem}
\begin{proof}
	لتكن \(B,C\) نقطتين مختلفتين على مستقيم \(l\) . نرسم مستقيمين عمودين من \(l\) على \(B,C\) وهذان المستقيمان سيتقاطعان في نقطة \(A\) .
	ولذلك فان مجموع زوايا المثلث اكبر من 180 .
\end{proof}

\begin{note}
	بين أي نقطتين على سطح كرة (ما عدا كونهما نهايتي قطر من أقطارها) يوجد خط واحد فقط، وهو جزء الدائرة الناتج من تقاطع المستوي المتعين بالنقطتين ومركز الكرة.
فمثلاً إذا كانت النقطتان على سطح الكرة هما:
$A(x_1, y_1, z_1)$ و $B(x_2, y_2, z_2)$
فإنهما تحققان معادلة الكرة، وتكون معادلة المستوي المار بالنقطتين $A$, $B$ ونقطة الأصل معطاة بالمحدد الآتي:

\[
\begin{vmatrix}
	x & y & z \\
	x_1 & y_1 & z_1 \\
	x_2 & y_2 & z_2
\end{vmatrix} = 0
\]

إن هذه المعادلة خطية من الدرجة الأولى بالنسبة للمتغيرات الثلاثة $x, y, z$، وتمر بالنقطة $(0,0,0)$ وتحقق النقطتين $A$ و $B$.

ولغرض دراسة الهندسة الإهليليجية على مستوى كما في الهندسة الهذلولية والإقليدية نستخدم الإسقاط الجسماني الذي ينقل جميع النقاط على سطح الكرة إلى المستوي $XY$، مع اعتبار النقطة $N(0,0,1)$ ثابتة (ليس لها مسقط أو أن مسقطها في مالانهاية).

فالإسقاط لأي نقطة $P(x,y,z)$ من سطح الكرة على المستوي $xy$ يتم بوصله بين النقطتين $P$ و $N$ ثم مدّه حتى يلاقي مستوي $XY$ في نقطة مثل $Q$، وتكون النقطة $Q$ هي مسقط النقطة $P$.

\begin{figure}[H]
	\raggedleft
	\includegraphics[scale=0.3]{plot11.pdf}
\end{figure}

إن هذا الإسقاط يؤلف إسقاطاً متبايناً عدا النقطة $N$، أي إن كل نقطة على سطح الكرة (ما عدا النقطة $N$) توجد لها نقطة واحدة فقط على المستوي $XY$ وبالعكس، كل نقطة على المستوي $XY$ توجد لها نقطة مقابلة واحدة على سطح الكرة.
\end{note}

\begin{example}
	 اوجد معادلة المستوي المار بنقطة الأصل والنقطتين
	 $A(1,2,3)$ و $B(4,5,6)$
\end{example}
\textbf{الحل}
\[
\begin{vmatrix}
	x & y & z \\
	x_1 & y_1 & z_1 \\
	x_2 & y_2 & z_2
\end{vmatrix} = 0
\Rightarrow 
\begin{vmatrix}
	x & y & z \\
	1 & 2 & 3 \\
	4 & 5 & 6
\end{vmatrix} = 0
\]
\[
\Rightarrow -3x+ 6y - 3z = 0 \Rightarrow x - 2y + z = 0
\]

\section[خواص الإسقاط الجسماني]{خواص الإسقاط الجسماني  \cite{ref6}}
\begin{enumerate}
	\item النقاط الواقعة على خط الاستواء \((z = 0)\) يكون إسقاطها على نفسها.
	\item النقاط الواقعة أعلى المستوى \(XY\) \((z > 0)\) من سطح الكرة تسقط خارج دائرة الاستواء.
	\item النقاط الواقعة أسفل المستوى \(XY\) \((z < 0)\) من سطح الكرة تسقط على نقاط داخل دائرة الاستواء.
\end{enumerate}

\textbf{ملاحظة:}
\begin{itemize}[label=\textbullet]
	\item إذا كانت النقطة \(Q(s,t)\) في المستوى \(XY\) هي مسقط النقطة \(P(x,y,z)\) الواقعة على سطح الكرة، فإنه يمكن إيجاد إحدى النقطتين بدلالة الأخرى.
	\item إذا عُلمت النقطة \(Q(s,t)\) فإن النقطة \(P(x,y,z)\) تُحسب بالعلاقات التالية:
\begin{align*}
	x &= \frac{2s}{s^2 + t^2 + 1},\\
	y &= \frac{2t}{s^2 + t^2 + 1},\\
	z &= \frac{s^2 + t^2 - 1}{s^2 + t^2 + 1}.
\end{align*}

	\item إذا عُلمت النقطة \(P(x,y,z)\) فإن النقطة \(Q(s,t)\) تُحسب بالعلاقات التالية:

\begin{align*}
	s &= \frac{x}{1 - z},\\
	t &= \frac{y}{1 - z}, \qquad -1 < z < 1.
\end{align*}
\end{itemize}

\begin{example}
	جد النقطة التي يكون مسقطها على $XY$ هو 
	$َQ(4, 6)$ على الكرة 
	$x^2 + y^2 + z^2 = 1$.
\end{example}
\begin{solution}
	من العلاقات:
	\[
	x = \frac{2s}{s^2 + t^2 + 1}, \quad
	y = \frac{2t}{s^2 + t^2 + 1}, \quad
	z = \frac{s^2 + t^2 - 1}{s^2 + t^2 + 1}
	\]
	إذن المقام:
	\[
	s^2 + t^2 + 1 = 52 + 1 = 53
	\]
	نحسب الإحداثيات:
	\[
	x = \frac{2 \cdot 4}{53} = \frac{8}{53}
	\]
	\[
	y = \frac{2 \cdot 6}{53} = \frac{12}{53}
	\]
	\[
	z = \frac{52 - 1}{53} = \frac{51}{53}
	\]
	إذن النقطة المطلوبة هي:
	\[
	P(x,y,z) = \left(\frac{8}{53}, \frac{12}{53}, \frac{51}{53}\right)
	\]
\end{solution}





\section[المسافة الإهليليجية بين نقطتين]{المسافة الإهليليجية بين نقطتين \cite{ref5}}

المسافة الإهليلجية بين النقطتين \(P_1(x_1,y_1,z_1)\) و \(P_2(x_2,y_2,z_2)\) على سطح الكرة تُمثَّل بالزاوية المركزية التي تقابل القوس الأصغر من الدائرة العظمى المارة بالنقطتين (ويجب أن تكون الزاوية مقاسة بالراديان). أي أن الزاوية \(\angle P_1 O P_2\) يرمز لها بـ \(\theta\)
\begin{figure}[H]
	\raggedleft
	\includegraphics[scale=0.3]{plot12.pdf}
\end{figure}
، وتحسب بالعلاقة:

\[
\cos \theta = \frac{(OP_1 \cdot OP_2)}{\|{OP_1}\| \, \|{OP_2}\|}
\]

وبما أن:
\[
\|{OP_1}\| = \|{OP_2}\| = 1
\]
(لأنها تمثل نصف قطر الكرة)

إذن:
\[
\cos \theta = x_1 x_2 + y_1 y_2 + z_1 z_2
\]

وإذا كان مسقط النقطتين الجسماني على المستوى \(XY\) هما:
\[
Q_1(s_1,t_1), \quad Q_2(s_2,t_2)
\]

فإن:
\[
\cos \theta =
\frac{
	4 s_1 s_2 + 4 t_1 t_2 + (s_1^2 + t_1^2 - 1)(s_2^2 + t_2^2 - 1)
}{
	(s_1^2 + t_1^2 + 1)(s_2^2 + t_2^2 + 1)
}
\]

\textbf{ملاحظات:}
\begin{enumerate}
	\item من القانون السابق نحسب قيمة \(\cos \theta\)، ثم نأخذ الدالة العكسية \(\arccos\) للحصول على \(\theta\) بوحدة الراديان، وهذه تمثل المسافة الإهليلجية بين \(P_1\) و\(P_2\).
	\item لتحويل الزاوية من الدرجات إلى الراديان نضربها في:
	\[
	\frac{\pi}{180}
	\]
\end{enumerate}

\begin{example}
	جد المسافة الإهليليجية بين النقطتين 
	$Q_1(6, 4)$ ، $Q_2(3, -7)$
\end{example}

\begin{solution}
	لدينا:
	\[
	Q_1(s_1,t_1) = (6,4), \quad Q_2(s_2,t_2) = (3,-7)
	\]
	نستخدم العلاقة:
	\[
	\cos \theta =
	\frac{
		4 s_1 s_2 + 4 t_1 t_2 + (s_1^2 + t_1^2 - 1)(s_2^2 + t_2^2 - 1)
	}{
		(s_1^2 + t_1^2 + 1)(s_2^2 + t_2^2 + 1)
	}
	\]
	نحسب أولاً:
	\[
	s_1^2 + t_1^2 = 36 + 16 = 52 \Rightarrow s_1^2 + t_1^2 + 1 = 53
	\]
	\[
	s_2^2 + t_2^2 = 9 + 49 = 58 \Rightarrow s_2^2 + t_2^2 + 1 = 59
	\]
	\[
	4 s_1 s_2 = 4(6)(3) = 72
	\]
	\[
	4 t_1 t_2 = 4(4)(-7) = -112
	\]
	\[
	(s_1^2 + t_1^2 - 1) = 51,\quad (s_2^2 + t_2^2 - 1) = 57
	\]
	
	إذن:
	\[
	\cos \theta = \frac{2867}{3127}
	\]
	ومن ثم:
	\[
	\theta = \arccos\left(\frac{2867}{3127}\right) \approx 23.5297
	\]
	بالتقدير الدائري
	\[
	\theta = 23.5297 \times \frac{\pi}{180} = 0.4016
	\]
\end{solution}
